\Needspace{9\baselineskip}
\originalchapter{稳定子、Clifford与拓扑码}
MAS.865J第4—6讲把纠错从逐码字计算改成Pauli代数. 其优势不是省略量子态, 而是发现某些态和门的全部可观测性质可以由一个多项式大小的对易表描述. 本章先证明维数与测量更新, 再从同一语言构造表面码, 说明简并错误为何不必逐个分辨.

\originalsection{Pauli串的二进制辛结构}
忽略整体相位, $n$位Pauli串可写$X^xZ^z$, $x,z\in\mathbb F_2^n$. 对任一位置$XZ=-ZX$, 两串交换产生的符号为
\[(-1)^{x\cdot z'+z\cdot x'}.\]
所以对易当且仅当这个辛内积为0. 相位不能永久丢掉: 稳定子要求本征值$+1$, 若两个生成元乘起来为$-I$, 它们的共同$+1$空间就是零. 二进制表还应保存符号位, 在乘法或Clifford共轭时按实际Pauli乘积更新.

稳定子群S取r个独立、两两对易的Hermitian Pauli生成元$g_j$, 不含$-I$. 共同$+1$投影
\[P=\prod_{j=1}^r\frac{I+g_j}2=2^{-r}\sum_{g\in S}g.\]
非恒等Pauli迹为0, 独立性使S中只有一项I, 因此$\Tr P=2^{n-r}$. 这既证明码空间维数, 也说明完整稳定子态$r=n$唯一. 减少一个独立检查会多留一个逻辑量子比特, 不能把重复生成元也计入r.

\begin{example}Bell态由$XX,ZZ$稳定, 两生成元独立, 维数$2^{2-2}=1$. 它们的乘积是$-YY$, 不是$+YY$. 三位重复码由$ZZI,IZZ$稳定, 维数2. 多加$XXX$会选出GHZ$+$, 多加$-XXX$选出GHZ$-$; 不加它则保留未知逻辑态.\end{example}

\originalsection{逻辑算符与简并错误的判别}
在Pauli群内, 与所有S对易的串构成中心化子. 由于Pauli共轭只会把一个稳定子变为$\pm g$, 而S不含$-I$, 群正规化子在这里也就是中心化子. 其中属于S的串在码上恒等; 不属于S的串则可能为非平凡逻辑Pauli.

用二进制矩阵看, r个独立辛约束留下$2n-r$维中心化子, 再除r维稳定子, 得$2(n-r)=2k$维逻辑Pauli空间. 可选k对辛反对易基$\overline X_j,\overline Z_j$, 各对像普通量子比特, 不同逻辑位的串对易. 选逻辑基必须在保存符号的实际Pauli表上核对, 不能只给一个无相位的二进制向量.

两错误$E_a,E_b$同综合征, 当且仅当$E_a^\dagger E_b$与S对易. 若乘积属于S至相位, 两者在码上相同, 同一个恢复即可, 是简并. 若乘积为非平凡逻辑串, 恢复无法知道哪种逻辑变化发生, KL条件失败. 若它与某个稳定子反对易, 两错误码像正交, 综合征能够区分. 因而量子距离是中心化子中不属于S的最小非零权重, 而不是所有稳定子的最小权重.

\originalsection{Pauli测量为何只需要更新表}
测Hermitian Pauli M, 投影为$P_m=(I+(-1)^mM)/2$. 若M与全部g对易, 对一个完整稳定子\textbf{纯态}, M或−M必在群中, 结果确定. 对只定义码空间的表, M还可能是一个逻辑测量, 会揭示逻辑信息; 这两情形不能混同.

若M与某个$g_1$反对易, 原态有$\langle M\rangle=\langle g_1Mg_1\rangle=-\langle M\rangle=0$, 所以两结果等概率. 把其余与M反对易的g先乘$g_1$, 它们就与M对易; 再将$g_1$换成$(-1)^mM$. 这些新检查固定投影后的态, 且仍独立, 因而完整描述后态. 测量成本就是Pauli乘法与符号更新, 无需存$2^n$个振幅.

\begin{example}对Bell表$XX,ZZ$测第二位Y. 两生成元都反对易$IY$. 留$g_2g_1=-YY$, 换第一为$\epsilon IY$, $\epsilon=\pm1$. 新表指明第二位Y本征值为$\epsilon$, 第一为$-\epsilon$, 后态是两个相反Y本征态的直积. 局部测量消耗了Bell纠缠, 而表计算保留了所有条件相关.\end{example}

\originalsection{Clifford门与Gottesman--Knill模拟}
Clifford门把Pauli共轭成Pauli. 单位H给$X\leftrightarrow Z,Y\mapsto-Y$; S给$X\mapsto Y,Y\mapsto-X,Z\mapsto Z$. CNOT控制c目标t给
\[X_c\mapsto X_cX_t,\quad Z_c\mapsto Z_c,\quad
X_t\mapsto X_t,\quad Z_t\mapsto Z_cZ_t.\]
对每个生成元逐位置应用规则即可. H、S的有向轴置换及符号旋转生成单比特24种Clifford, 只差整体相位; 多位用CNOT把串支持消去, 可把一般Clifford分解为$O(n^2)$个H、S、CNOT, PS2 P4给实际归纳块分解.

若初态是稳定子态, 电路只有Clifford和Pauli测量及按结果选Clifford, 就可维护$O(n^2)$大小的表. 每个门和测量更新多项式, 抽随机测量标签后继续更新, 精确模拟整个分布. 这是Gottesman--Knill的构造证明, 不是声称“有纠缠也完全没用”. Bell、GHZ甚至大量纠缠态在这个受限门族中都可模拟. T门把X变成$(X+Y)/\sqrt2$, 超出单Pauli表, 是普适门库中破坏这一闭合的一项.

\originalsection{横向门怎样防传播}
编码门横向连接不同块的同一物理位置, 一个两位门故障至多在每块产生一个位置错误. CNOT的共轭表也给传播方向: X从控制传目标, Z从目标传控制. 在同一块内让一个辅助依次连许多数据, 就可能把一错传播成高权重错误. 这解释为何“逻辑门保码”与“逻辑门容错”是两条不同证明.

Steane横向H通过第6章的码字Fourier计算确为逻辑H. 横向CNOT可在两个CSS陪集和上逐基计算: 目标标签$x_2+y_2$异或控制$x_1+y_1$, y的线性空间重排不变, 实现逻辑异或. 对S则要检查码字重量相位, 对一般码没有统一横向结论; Bacon--Shor的S注入与验证将在第8章实际构造.

\originalsection{环面码的检查、维数与逻辑环}
在周期$L\times L$方格每条边放一个量子比特, 共$E=2L^2$位. 每个顶点星$A_v$为相邻四边X积, 每个面$B_f$为边界四边Z积. 星与面相交0或2边, 辛内积偶, 故全部对易. 全部星的乘积I, 因每边出现两次; 全部面的乘积也I. 除这两关系没有额外关系, 独立检查$(V-1)+(F-1)=2L^2-2$, 所以编码2位.

Z错误串的端点会使相应星的本征值翻转. 闭合Z环没有端点, 不改变综合征. 若环可收缩, 它是所包围面的$B_f$乘积, 对码态恒等; 若绕环面一周不可收缩, 它是逻辑Z. 对偶图的X串同理给逻辑X. 一条X逻辑环与对应Z逻辑环交一边时反对易, 不同方向的适当环可对易, 因而产生两对逻辑Pauli. 最短非平凡环有L边, 码距L.

从综合征恢复时, 可以任选路径连接相同端点. 两条恢复路径相差一个闭环; 只要闭环可收缩, 两恢复只差稳定子, 都成功. 这就是大量物理路径不必分别定位的简并纠错. 若真实错误与恢复相差不可收缩环, 会发生未被检测的逻辑错, 所以解码的统计权重与拓扑类很重要.

\originalsection{曲面、边界与任意子解释}
一般封闭曲面仍有$V+F-2$个独立检查, 编码$E-V-F+2$. 环面Euler数0给k=2; 射影平面Euler数1在模2胞腔结构下给k=1. 对有重复边的面界, 必须按模2出现次数计数: 边出现两次抵消. 原作业只用直觉图规则声称重复无效并不普遍正确, 题解给具体CW反例与合法细分.

将星或面−1的激发叫电、磁任意子, 是二维模型中的标签. 用Z串搬运一类激发, X串搬运另一类; 两串交一次时XZ反对易, 相应绕行过程相差−1相位. 这给互统计相位, 不能据此把普通三维粒子的任意交换都叫任意子. 非Abel任意子的融合空间可以承载量子信息, 但本章建立的是Abel表面码及其检查语言, 不把非Abel通用容错方案当已证明.
